\chapter{積分と数列}\label{章:積分と数列}

項数の増えていく和は,そのままでは計算しにくい.しかし,和を長方形の面積の和と思えば,積分の値として求められることがある.また,積分の中身を大小比較すれば,値を正確に求められなくても積分の範囲が分かる.この2つを組み合わせると,数列の極限や,複雑な積の極限まで調べられる.

この章では,まず和の中に区分求積の形を見つける.その後,積分を評価する不等式と,指数を1つずつ変えた積分を結ぶ漸化式を考える.漸化式は部分積分から得られ,数列の漸化式と同じ手順で解ける.最後のベータ積分では,それまで別々に扱った多項式の積分と三角関数の積分を,1つの式の特別な場合として並べる.

\section{区分求積法}\label{節:和の区分求積}

次の和は,\,$n$を大きくするとどうなるだろうか.
\[
 S_n=\sum_{k=1}^{n}\frac{n}{n^2+k^2}
\]
各項の分母は$n^2$以上なので,どの項も$\frac1n$以下であり,\,$n\to\infty$で$0$へ近づく.一方で,足す項の個数$n$は増えていく.どちらの影響が強いのかは,式を眺めるだけでは分からない.そこで実際に計算してみよう.
\begin{center}
\begin{tabular}{ccccccc}
\toprule
$n$ & $1$ & $2$ & $5$ & $10$ & $100$ & $1000$\\
\midrule
$S_n$ & $0.5$ & $0.65$ & $0.7337$ & $0.7600$ & $0.7829$ & $0.7851$\\
\bottomrule
\end{tabular}
\end{center}
$S_n$は$0$へは向かわず,\,$0.785$付近へ近づいているように見える.この値の正体を突き止めよう.

ここで1つの知識を思い出そう.定積分は,区間を細かく分けて作った長方形の面積の和の極限であった.そこで,和の極限を定積分に帰着させることを目標とする.

連続関数$f$について,区間$[a,b]$を$n$等分し,各小区間の右端での関数値と幅$\frac{b-a}{n}$を掛けて足すと,
\[
 \frac{b-a}{n}\sum_{k=1}^{n}
 f\left(a+\frac{k(b-a)}{n}\right)
\]
となる.\,$f$が0以上なら長方形の面積の和であり,負になる部分は符号を付けて足している.\,$n$を大きくすれば,この和は$\int_a^b f(x)\dd x$へ近づく.とくに$[0,1]$では
\begin{equation}\label{式:単位区間の区分求積}
 \lim_{n\to\infty}\frac1n\sum_{k=1}^{n}f\left(\frac{k}{n}\right)
 =\int_0^1 f(x)\dd x
\end{equation}
となる.この式を使って和の極限を求める方法を\textbf{区分求積法}という.

使うときには,\,$\frac kn$が横の位置,\,$\frac1n$が長方形の幅であることを確かめる.分母に$n$が現れるというだけでは足りない.和全体を,幅と高さの積の和へ書き直す必要がある.

\begin{bookpoint}
\textbf{POINT}\par
和から幅$\frac1n$をくくり出し,残りを$\frac kn$だけの式$f\left(\frac kn\right)$に書き直して,\,$\frac kn$が動く範囲で$f$を積分する.
\end{bookpoint}

\begin{example}{和から積分へ}
次の極限を求めよ.
\[
 \lim_{n\to\infty}\sum_{k=1}^{n}\frac{n}{n^2+k^2}.
\]
\end{example}

\paragraphbf{方針}
各項の分母から$n^2$をくくり,幅$\frac1n$と$\frac kn$の関数の積に書き直す.

\paragraphbf{解答}
$\frac{\pi}{4}$

\paragraphbf{解法}
各項の分母から$n^2$をくくると
\[
 \frac{n}{n^2+k^2}
 =\frac1n\cdot\cfrac{1}{1+\left(\cfrac{k}{n}\right)^2}
\]
となる.したがって,これは$f(x)=\frac1{1+x^2}$の右端を使う区分求積である.この関数は$[0,1]$で連続なので
\[
 \lim_{n\to\infty}\sum_{k=1}^{n}\frac{n}{n^2+k^2}
 =\int_0^1\frac{1}{1+x^2}\dd x
\]
を得る.右辺で$x=\tan t$と置くと,\,$x$が$0$から$1$まで動くとき$t$は$0$から$\frac\pi4$まで動き,\,$\dd x=(1+\tan^2t)\dd t$である.よって
\[
 \int_0^1\frac{1}{1+x^2}\dd x
 =\int_0^{\frac\pi4}\frac{1+\tan^2t}{1+\tan^2t}\dd t
 =\frac{\pi}{4}.
\]
表で見た$0.785$付近の値は,\,$\frac\pi4=0.78539\ldots$だったのだ.

\paragraphbf{積分区間の決定}
和に含まれる点の範囲が変わるなら,積分区間も確かめなければならない.例えば
\[
 \frac1n\sum_{k=1}^{2n}f\left(\frac{k}{n}\right)
\]
では,幅は$\frac1n$のままで,点は$\frac1n$から$2$まで並ぶ.これは$[0,2]$の区分求積である.項数が$2n$だからといって,幅を$\frac1{2n}$に変えてはいけない.

\begin{example}{積分区間の決定}
次の極限を求めよ.
\[
 \lim_{n\to\infty}\sum_{k=n+1}^{2n}\frac1k.
\]
\end{example}

\paragraphbf{方針}
各項を$\frac1n\cdot\cfrac1{\cfrac kn}$と書き,\,$k=n+1,\ldots,2n$に対応する$\frac kn$の範囲から積分区間を決める.

\paragraphbf{解答}
$\log2$

\paragraphbf{解法}
各項を
\[
 \frac1k=\frac1n\cdot\cfrac1{\cfrac{k}{n}}
\]
と書く.\,$k=n+1,\ldots,2n$に対応する点$\frac kn$は$1+\frac1n,\ldots,2$であり,\,$[1,2]$を$n$等分した小区間の右端である.\,$\frac1x$は$[1,2]$で連続なので
\[
 \lim_{n\to\infty}\sum_{k=n+1}^{2n}\frac1k
 =\int_1^2\frac1x\dd x
 =\log2.
\]
$\frac1x$は$x=0$で定義されないが,使った積分区間は$[1,2]$なので問題にならない.

\paragraphbf{端の項の扱い}
式\eqref{式:単位区間の区分求積}で,右端の代わりに左端や中点を使っても,同じ積分へ収束する.左端を使う和は,右端を使う和と端の1項ずつが違うだけである.実際,連続関数$f$について
\[
 \frac1n\sum_{k=0}^{n}f\left(\frac{k}{n}\right)
 -\frac1n\sum_{k=1}^{n}f\left(\frac{k}{n}\right)
 =\frac{f(0)}n
 \longrightarrow0
\]
である.このように,有限個の端の項を付けたり取り除いたりしても,その差が$0$へ近づくなら極限は変わらない.しかし,取り除いた項数が$n$とともに増える場合や,端の関数値が大きくなる場合には,この確認を省けない.区分求積の公式を使う範囲を先に決め,残りの差は別に評価すると良い.
%
\section{和と積分の比較}\label{節:和と積分の比較}

区分求積では,極限を取って和と積分を結び付けた.単調性が分かる場合には,有限の$n$でも大小関係が得られる.\,$f$が$[1,\infty)$で正かつ単調減少なら,\,$k\leqq x\leqq k+1$において
\[
 f(k+1)\leqq f(x)\leqq f(k)
\]
である.長さ$1$の区間で積分すると
\[
 f(k+1)\leqq\int_k^{k+1}f(x)\dd x\leqq f(k).
\]
この不等式を足すと
\begin{equation}\label{式:減少関数の和比較}
 \int_1^{n+1}f(x)\dd x
 \leqq\sum_{k=1}^{n}f(k)
 \leqq f(1)+\int_1^{n}f(x)\dd x
\end{equation}
を得る.右側では,\,$k=2,\ldots,n$について$f(k)\leqq\int_{k-1}^{k}f(x)\dd x$を足し,最初の$f(1)$だけを別に残している.

図\ref{図:和と積分の比較}では,各長方形の高さを小区間の左端の値にしている.減少関数では長方形が曲線の上に出るので,長方形の面積の和が積分の上からの評価になる.高さを右端の値に変えれば,大小関係が逆になる.

\begin{figure}[htbp]
 \centering
 \begin{tikzpicture}[x=1.15cm,y=2.9cm,>=Stealth]
  \foreach \k in {1,2,3,4} {
   \pgfmathsetmacro{\h}{1/\k}
   \fill[blue!12] (\k,0) rectangle ({\k+1},\h);
   \draw[blue!65!black] (\k,0) rectangle ({\k+1},\h);
  }
  \draw[->] (0.55,0)--(5.45,0) node[right]{$x$};
  \draw[->] (0.7,0)--(0.7,1.25) node[above]{$y$};
  \draw[thick,domain=1:5,samples=90] plot (\x,{1/\x});
  \node[anchor=west] at (2,0.86) {$y=\frac1x$};
  \foreach \k in {1,2,3,4,5}
   \node[below] at (\k,0) {$\k$};
 \end{tikzpicture}
 \caption{減少関数の左端の値で作った長方形は,曲線の上に出る.}
 \label{図:和と積分の比較}
\end{figure}

\begin{example}{調和数}
$H_n=\sum_{k=1}^{n}\frac1k$とおく.\,$H_n\to\infty$であること,および
\[
 \lim_{n\to\infty}\frac{H_n}{\log n}=1
\]
を示せ.後の比は$n\geqq2$について考える.
\end{example}

式\eqref{式:減少関数の和比較}を$f(x)=\frac1x$に用いると
\[
 \log(n+1)\leqq H_n\leqq1+\log n
\]
となる.下側の$\log(n+1)$が限りなく大きくなるため,\,$H_n$も限りなく大きくなる.次に,正の$\log n$で割れば
\[
 \frac{\log(n+1)}{\log n}
 \leqq\frac{H_n}{\log n}
 \leqq1+\frac1{\log n}.
\]
左辺は
\[
 \frac{\log(n+1)}{\log n}
 =1+\cfrac{\log\left(1+\cfrac1n\right)}{\log n}
 \longrightarrow1
\]
であり,右辺も$1$へ近づく.はさみうちの原理から求める極限が得られる.

ここから$H_n-\log n\to0$と結論してはいけない.比が$1$へ近づくということは,差が元の大きさに比べて小さいことを表すだけである.実際,上の評価から分かるのは,差が$0$以上$1$以下の範囲にあることまでである.

比較の向きを図から毎回考える必要はない.増加関数なら$f(k)\leqq f(x)\leqq f(k+1)$,減少関数ならその逆,と小区間ごとの不等式を書けばよい.添字のずれを避けるためにも,最初から大きな和の不等式を暗記するより,小区間から組み立てる方が扱いやすい.


%
\section{積分の不等式}\label{節:積分の不等式}

$f(x)\le g(x)$が区間全体で成り立てば,\,$\int_a^b f(x)\dd x\le\int_a^b g(x)\dd x$である.この基本的な性質を積や平方の形に使うと,2つの関数の積分の間に成り立つ不等式が得られる.

\begin{theorem}{積分に対するコーシー・シュワルツの不等式}
$a<b$とし,\,$f,g$は$[a,b]$で実数値を取る連続関数とする.このとき
\begin{equation}\label{式:積分のコーシーシュワルツ}
 \left(\int_a^b f(x)g(x)\dd x\right)^2
 \le
 \left(\int_a^b f(x)^2\dd x\right)
 \left(\int_a^b g(x)^2\dd x\right)
\end{equation}
が成り立つ.\,$g$が恒等的に$0$でなければ,等号が成り立つための必要十分条件は,ある定数$c$について区間全体で$f(x)=cg(x)$となることである.\,$g$が恒等的に$0$なら,どの$f$についても等号が成り立つ.
\end{theorem}

平方は負にならないので,任意の実数$t$について
\[
 0\le\int_a^b\{f(x)-tg(x)\}^2\dd x
\]
である.式を短くするため
\[
 A=\int_a^b f(x)^2\dd x,\qquad
 B=\int_a^b f(x)g(x)\dd x,\qquad
 C=\int_a^b g(x)^2\dd x
\]
と置くと,右辺を展開して
\[
 0\le Ct^2-2Bt+A
\]
となる.

$C>0$とする\footnote{$C=0$なら,下の等号の議論と同じく連続性から$g$は区間全体で$0$となり,\,$B=0$である.このとき不等式の両辺はともに$0$である.}.右辺は$t$の2次関数であり,すべての実数$t$で$0$以上である.ここで1つの知識を思い出そう.2次関数がすべての実数で$0$以上なら,その判別式は$0$以下である.右辺の判別式を$D$とすると
\[
 \frac{D}{4}=B^2-AC\le0
\]
となり,\,$B^2\le AC$を得る.これが式\eqref{式:積分のコーシーシュワルツ}である.

等号$B^2=AC$が成り立つのは判別式が$0$のときであり,そのとき2次関数は$t=\frac BC$で$0$になる.つまり
\[
 \int_a^b\left\{f(x)-\frac BCg(x)\right\}^2\dd x=0.
\]
被積分関数は連続かつ$0$以上である.もし1点で正なら,連続性によってその近くの小区間でも正となり,積分が正になってしまう.したがって区間全体で$f(x)=\frac BCg(x)$である.逆に$f=cg$なら,直接代入して等号を確かめられる.

この不等式を使うときには,評価したい積分の中身を$f(x)g(x)$という積に分ける.片方を$1$とすることも多い.例えば$g(x)=1$とすれば
\begin{equation}\label{式:平均と二乗平均}
 \left(\int_a^b f(x)\dd x\right)^2
 \le(b-a)\int_a^b f(x)^2\dd x
\end{equation}
となる.両辺を$(b-a)^2$で割ると,関数の平均値の平方は,平方の平均値を超えないことが分かる.等号は$f$が定数関数のときに限る.

\begin{bookpoint}
\textbf{POINT}\par
評価したい積分の中身を$f(x)g(x)$の積に分け,\,$f^2$と$g^2$の積分が求めたい形になるように$f,g$を選んで式\eqref{式:積分のコーシーシュワルツ}に当てはめる.
\end{bookpoint}

\begin{example}{逆数と不等式}
$h$は$[a,b]$で連続で,すべての$x$について$h(x)>0$を満たすとする.
\[
 \left(\int_a^b h(x)\dd x\right)
 \left(\int_a^b\frac1{h(x)}\dd x\right)
 \ge(b-a)^2
\]
を示し,等号条件を求めよ.
\end{example}

\paragraphbf{方針}
2つの積分を右辺の$\int f^2$と$\int g^2$にするため,\,$f=\sqrt h$,\,$g=\frac1{\sqrt h}$と置く.積$fg$は$1$になる.

\paragraphbf{解答}
不等式は成り立ち,等号は$h$が正の定数関数のとき,またそのときに限り成り立つ.

\paragraphbf{解法}
$f=\sqrt h$,\,$g=\frac1{\sqrt h}$と置く.\,$h$は閉区間で連続かつ正なので,この2つの関数も連続である.\,$f^2=h$,\,$g^2=\frac1h$,\,$fg=1$だから,式\eqref{式:積分のコーシーシュワルツ}より
\[
 (b-a)^2
 =\left(\int_a^b1\dd x\right)^2
 \le\left(\int_a^b h(x)\dd x\right)
       \left(\int_a^b\frac1{h(x)}\dd x\right)
\]
となる.\,$g$は恒等的に$0$ではないので,等号条件は$f=cg$である.これは$\sqrt h=\frac c{\sqrt h}$,すなわち$h=c$を意味する.

\paragraphbf{対数の評価}
例えば$h(x)=1+x$,\,$0\le x\le1$とすると,
\[
 \int_0^1(1+x)\dd x=\frac32,\qquad
 \int_0^1\frac1{1+x}\dd x=\log2
\]
である.\,$h$は定数ではないから等号は成り立たず,
\[
 \frac32\log2>1
\]
となる.つまり$\log2>\frac23=0.666\ldots$である.実際$\log2=0.693\ldots$であり,対数の値を小数で計算せずに,厳密な下からの評価が得られたのだ.
%
\section{積分と漸化式}\label{節:積分漸化式}

積分に自然数$n$が含まれているとき,\,$n=0,1,2,\ldots$の値を1つずつ計算するのは手間である.例えば$\int_0^1x^n\e^{-x}\dd x$を直接求めるには,部分積分を$n$回繰り返さなければならない.そこで,\,$n$を1つ小さくした積分との関係,つまり漸化式を作ることを目標とする.

手がかりは部分積分である.\,$x^n$を微分すると$nx^{n-1}$となり,指数が1つ下がる.相手の関数$g(x)$の原始関数を$G(x)$とすると,\,$n\ge1$で
\[
 \int_a^b x^ng(x)\dd x
 =\left[x^nG(x)\right]_a^b
 -n\int_a^b x^{n-1}G(x)\dd x
\]
となる.\,$g(x)=\e^{-x}$のように,\,$G(x)$が$g(x)$の定数倍になる関数なら,右辺の積分は左辺の$n$を$n-1$に変えた積分の定数倍である.こうして隣り合う$n$の積分が漸化式で結ばれる.

\begin{bookpoint}
\textbf{POINT}\par
$x^n$を微分する側に選んで部分積分し,\,$n$番目の積分を$n-1$番目の積分で表す.最初の値は直接計算する.
\end{bookpoint}

\begin{example}{積分と漸化式}
0以上の整数$n$に対して
\[
 A_n=\int_0^1 x^n\e^{-x}\dd x
\]
とおく.\,$A_n$の漸化式を求め,\,$A_2$を計算せよ.また$A_n\to0$を示せ.
\end{example}

\paragraphbf{方針}
$x^n$を微分し,\,$\e^{-x}$を積分する部分積分で$A_n$を$A_{n-1}$で表す.極限は漸化式を使わず,\,$\e^{-x}\le1$で積分を上から抑える.

\paragraphbf{解答}
$A_n=nA_{n-1}-\e^{-1}$\,$(n\ge1)$,\,$A_0=1-\e^{-1}$,\,$A_2=2-5\e^{-1}$.\,$A_n\to0$も成り立つ.

\paragraphbf{解法}
$n\ge1$のとき,\,$x^n$を微分し,\,$\e^{-x}$を積分する部分積分を行うと
\begin{align*}
 A_n
 &=\left[-x^n\e^{-x}\right]_0^1
   +n\int_0^1x^{n-1}\e^{-x}\dd x\\
 &=-\e^{-1}+nA_{n-1}
\end{align*}
となる.\,$x=0$の端点の項が$0$になるのは$n\ge1$だからである.最初の値は
\[
 A_0=\int_0^1\e^{-x}\dd x=1-\e^{-1}
\]
と直接求める.漸化式より
\begin{align*}
 A_1&=1\cdot A_0-\e^{-1}=1-2\e^{-1},\\
 A_2&=2A_1-\e^{-1}=2-5\e^{-1}
\end{align*}
を得る.

極限を調べるには,漸化式より積分の評価が短い.\,$[0,1]$では$0<\e^{-x}\le1$なので
\[
 0<A_n\le\int_0^1x^n\dd x=\frac1{n+1}
\]
である.右辺は$0$へ近づくので,はさみうちの原理より$A_n\to0$である.

\paragraphbf{漸化式を解く}
得られた漸化式
\[
 A_n=nA_{n-1}-\e^{-1}
\]
を一般項まで解いてみよう.\,$A_{n-1}$に掛かる倍率$n$が添字とともに変わるので,特性方程式型のように定数を引いて等比型にすることはできない.そこで,倍率$n,n-1,\ldots,1$をまとめた$n!$で両辺を割ってやる.
\[
 \frac{A_n}{n!}=\frac{A_{n-1}}{(n-1)!}-\frac{\e^{-1}}{n!}
\]
ここで$B_n=\frac{A_n}{n!}$と置くと
\[
 B_n=B_{n-1}-\frac{\e^{-1}}{n!}
\]
となり,隣り合う項の差が$n$の式で与えられる階差型に帰着する.\,$k=1,\ldots,n$について足せば途中の項が消え,
\[
 B_n=B_0-\e^{-1}\sum_{k=1}^{n}\frac1{k!}
\]
を得る.\,$B_0=A_0=1-\e^{-1}$であり,\,$0!=1$なので
\[
 A_n=n!\left(1-\e^{-1}\sum_{k=0}^{n}\frac1{k!}\right)
\]
となる.ただし,この形から$n\to\infty$の極限を求めるには,括弧の中が$0$へ近づく速さと$n!$が大きくなる速さを比べなければならない.先ほどの積分による評価は,その比較をせずに直接$A_n\to0$を示している.

\paragraphbf{漸化式で値を求める実験}
漸化式があれば,\,$A_0$から順に値を計算できる.そこで$\e^{-1}=0.367879\ldots$を小数第4位までに丸めた$0.3679$を使い,\,$A_0=1-0.3679$から計算してみよう.比べるのは,積分の値を小数第4位まで求めたものである.
\begin{center}
\begin{tabular}{crr}
\toprule
$n$ & 積分の値 & 漸化式で求めた値\\
\midrule
$0$ & $0.6321$ & $0.6321$\\
$1$ & $0.2642$ & $0.2642$\\
$2$ & $0.1606$ & $0.1605$\\
$3$ & $0.1139$ & $0.1136$\\
$4$ & $0.0878$ & $0.0865$\\
$5$ & $0.0713$ & $0.0646$\\
$6$ & $0.0599$ & $0.0197$\\
$7$ & $0.0517$ & $-0.2300$\\
$8$ & $0.0454$ & $-2.2079$\\
\bottomrule
\end{tabular}
\end{center}
被積分関数は正なので$A_n>0$のはずだが,漸化式で求めた値は$n=7$で負になってしまった.

原因は丸めたことによるずれ$\delta=0.3679-\e^{-1}$である.漸化式で求めた値を$a_n$とし,\,$D_n=a_n-A_n$と置く.2つの漸化式
\[
 \begin{array}{rrl}
 &a_n={}&na_{n-1}-0.3679\\
 -\,)&A_n={}&nA_{n-1}-\e^{-1}\\ \hline
 &D_n={}&nD_{n-1}-\delta
 \end{array}
\]
の引き算から,\,$D_n$も同じ形の漸化式を満たす.\,$D_0=-\delta$なので,先ほどと同じく$n!$で割って足せば
\[
 D_n=-\delta\,n!\sum_{k=0}^{n}\frac1{k!}
\]
となる.\,$\delta$は$0.00002$程度にすぎないが,\,$n!$倍されて現れたのだ.\,$n=7$では$7!=5040$であり,ずれは約$0.28$になる.漸化式は値を順に求めるのに役立ち,積分表示は符号や大きさを調べるのに役立つ.同じ数列でも,問いに応じて使いやすい表し方が違う.

\paragraphbf{ネイピア数の級数}
一般項と積分による評価を組み合わせると,ネイピア数$\e$を和で表せる.一般項の両辺を$n!$で割って$\e$を掛けると
\[
 \frac{\e A_n}{n!}=\e-\sum_{k=0}^{n}\frac1{k!}
\]
となる.\,$0<A_n\le\frac1{n+1}$より
\[
 0<\e-\sum_{k=0}^{n}\frac1{k!}\le\frac{\e}{(n+1)!}
\]
である.右辺は$0$へ近づくので
\[
 \lim_{n\to\infty}\sum_{k=0}^{n}\frac1{k!}=\e
\]
を得る.例えば$n=5$なら和は$2.71666\ldots$であり,\,$\e=2.71828\ldots$との差は$\frac{\e}{6!}<0.004$以内に収まっている.
%
\section{ウォリスの公式}\label{節:ウォリス公式}

前節の方法を三角関数のべきに使ってみよう.
\[
 I_n=\int_0^{\frac{\pi}{2}}\sin^n x\dd x
 \qquad(n=0,1,2,\ldots)
\]
を考える.最初の2つは
\[
 I_0=\frac{\pi}{2},\qquad I_1=1
\]
である.\,$n\geqq2$の$I_n$を,指数の小さい積分に帰着させることを目標とする.

$\sin^n x$を$\sin^{n-1}x\cdot\sin x$と分け,後ろの$\sin x$を積分する側にして部分積分する.微分された$\sin^{n-1}x$から$\cos x$が現れ,積分の中に$\cos^2x$ができる.これを$1-\sin^2x$に書き直せば,再び$\sin$のべきだけの積分になる.

ここで1つ断りをいれておこう.\,$n\geqq2$なので$\sin^{n-1}0=0$であり,また$\cos\frac\pi2=0$である.したがって部分積分の端点の項は両端とも$0$になる.
\begin{align*}
 I_n
 &=\left[-\sin^{n-1}x\cos x\right]_0^{\frac{\pi}{2}}
 +(n-1)\int_0^{\frac{\pi}{2}}\sin^{n-2}x\cos^2x\dd x\\
 &=(n-1)\int_0^{\frac{\pi}{2}}\sin^{n-2}x(1-\sin^2x)\dd x\\
 &=(n-1)(I_{n-2}-I_n)
\end{align*}
最後の式の$I_n$を左辺に移項して$n$で割ると
\begin{equation}\label{式:ウォリス漸化式}
 I_n=\frac{n-1}{n}I_{n-2}\qquad(n\geqq2)
\end{equation}
を得る.

この漸化式は,\,$I_n$を2つ前の$I_{n-2}$と結ぶ.指数が$2$ずつ下がるので,偶数番目は$I_0$へ,奇数番目は$I_1$へつながる.偶数番目だけを取り出して$u_m=I_{2m}$と置くと,\,$u_m=\frac{2m-1}{2m}u_{m-1}$となり,隣り合う項の比が$m$の式で与えられる階比型になる.よって比を順に掛けてやれば
\begin{align*}
 I_{2m}
 &=\frac{2m-1}{2m}I_{2m-2}\\
 &=\frac{2m-1}{2m}\cdot\frac{2m-3}{2m-2}I_{2m-4}\\
 &\qquad\vdots\\
 &=\frac{2m-1}{2m}\cdot\frac{2m-3}{2m-2}\cdots\frac12I_0
\end{align*}
となる.奇数番目も同じく$I_1$まで下げればよい.

\begin{theorem}{ウォリスの積分公式}
$m\geqq1$に対して
\begin{equation}\label{式:ウォリス偶数}
 I_{2m}
 =\frac{1\cdot3\cdot5\cdots(2m-1)}{2\cdot4\cdot6\cdots2m}
   \cdot\frac{\pi}{2},
\end{equation}
\begin{equation}\label{式:ウォリス奇数}
 I_{2m+1}
 =\frac{2\cdot4\cdot6\cdots2m}{3\cdot5\cdot7\cdots(2m+1)}.
\end{equation}
\end{theorem}

偶数の積と奇数の積を短く書くために
\[
 (2m)!!=2\cdot4\cdots2m,\qquad
 (2m-1)!!=1\cdot3\cdots(2m-1)
\]
と定義することがある.これは\textbf{二重階乗}と呼ばれ,階乗を2回行うという意味ではない.この記号なら
\[
 I_{2m}=\frac{(2m-1)!!}{(2m)!!}\cdot\frac{\pi}{2},\qquad
 I_{2m+1}=\frac{(2m)!!}{(2m+1)!!}
\]
となる.

\begin{bookpoint}
\textbf{POINT}\par
$\sin^nx=\sin^{n-1}x\cdot\sin x$と分けて部分積分し,\,$\cos^2x=1-\sin^2x$で$I_n$を戻して移項する.\,$I_n=\frac{n-1}{n}I_{n-2}$を$I_0$か$I_1$まで繰り返す.
\end{bookpoint}

\begin{example}{ウォリス積分}
次の2つの積分を求めよ.
\[
 \int_0^{\frac{\pi}{2}}\sin^6x\dd x,\qquad
 \int_0^{\frac{\pi}{2}}\sin^7x\dd x.
\]
\end{example}

\paragraphbf{方針}
漸化式\eqref{式:ウォリス漸化式}で,偶数は$I_0=\frac\pi2$まで,奇数は$I_1=1$まで指数を下げる.

\paragraphbf{解答}
$\int_0^{\frac{\pi}{2}}\sin^6x\dd x=\frac{5\pi}{32}$,\qquad
$\int_0^{\frac{\pi}{2}}\sin^7x\dd x=\frac{16}{35}$

\paragraphbf{解法}
漸化式\eqref{式:ウォリス漸化式}を繰り返すと
\begin{align*}
 I_6&=\frac56\cdot\frac34\cdot\frac12\cdot\frac{\pi}{2}
     =\frac{5\pi}{32},\\
 I_7&=\frac67\cdot\frac45\cdot\frac23\cdot1
     =\frac{16}{35}
\end{align*}
となる.偶数と奇数で最後に戻る値が違うことが,答えに$\pi$が現れるかどうかの違いになっている.

積分区間が$\left[0,\frac{\pi}{2}\right]$であることも大切である.任意の区間では,部分積分の端点の項が同じように消えるとは限らない.

\subsection{積分の比}

公式は積分の値を積で表したものである.今度は逆に,積分の大小を使ってその積の極限を求める.まず$I_n$の値を並べてみよう.
\begin{center}
\begin{tabular}{cccccccccc}
\toprule
$n$ & $0$ & $1$ & $2$ & $3$ & $4$ & $5$ & $6$ & $7$ & $8$\\
\midrule
$I_n$ & $1.5708$ & $1$ & $0.7854$ & $0.6667$ & $0.5890$ & $0.5333$ & $0.4909$ & $0.4571$ & $0.4295$\\
\bottomrule
\end{tabular}
\end{center}
$I_n$は減少し,しかも隣り合う値の差が次第に小さくなっていく.隣どうしの比$\frac{I_{n+1}}{I_n}$は$1$へ近づくだろう.

実際,\,$0\le\sin x\le1$だから,指数を大きくすると被積分関数は小さくなる.したがって
\[
 I_{n+2}\le I_{n+1}\le I_n
\]
である.各辺を$I_n>0$で割り,式\eqref{式:ウォリス漸化式}を使うと
\[
 \frac{n+1}{n+2}
 =\frac{I_{n+2}}{I_n}
 \le\frac{I_{n+1}}{I_n}\le1
\]
となる.両端が$1$へ近づくため
\begin{equation}\label{式:ウォリス隣接比}
 \lim_{n\to\infty}\frac{I_{n+1}}{I_n}=1
\end{equation}
である.各$I_n$の極限を先に求めなくても,比の極限は分かるのだ.

ここで,次の有限の積を作る.
\[
 P_m=\prod_{k=1}^{m}\frac{(2k)^2}{(2k-1)(2k+1)}
\]
$\prod$は,指定された範囲の項をすべて掛ける記号である.最初の数項を書けば
\[
 P_m=\frac{2\cdot2}{1\cdot3}\cdot
      \frac{4\cdot4}{3\cdot5}\cdots
      \frac{2m\cdot2m}{(2m-1)(2m+1)}
\]
となる.分子の偶数の積は$\{(2m)!!\}^2$,分母の奇数の積は$(2m-1)!!\,(2m+1)!!$である.一方,式\eqref{式:ウォリス偶数}と式\eqref{式:ウォリス奇数}の比を取ると
\begin{align*}
 \frac{I_{2m+1}}{I_{2m}}
 &=\frac{2}{\pi}\cdot
   \frac{\{(2m)!!\}^2}{(2m-1)!!\,(2m+1)!!}\\
 &=\frac{2}{\pi}P_m
\end{align*}
となる.左辺は$1$へ近づくので,次の極限が得られる.
\begin{equation}\label{式:ウォリスの積}
 \boxed{
 \lim_{m\to\infty}\prod_{k=1}^{m}
 \frac{4k^2}{4k^2-1}=\frac{\pi}{2}}
\end{equation}
これが\textbf{ウォリスの積}である.無限個の数を一度に掛けたと考えるのではなく,有限の積$P_m$の数列の極限として定義している.

\paragraphbf{円周率を挟む}
証明の途中の不等式を残しておけば,近似の誤差も評価できる.\,$n=2m$とすると
\[
 \frac{2m+1}{2m+2}
 \le\frac{I_{2m+1}}{I_{2m}}\le1
\]
だから
\[
 \frac{2m+1}{2m+2}\cdot\frac{\pi}{2}
 \le P_m\le\frac{\pi}{2}
\]
である.各辺に正の数を掛けて$\pi$について解くと
\[
 2P_m\le\pi\le\frac{2m+2}{2m+1}\cdot2P_m
\]
となる.実際に計算してみよう.
\begin{center}
\begin{tabular}{ccc}
\toprule
$m$ & $2P_m$ & $\frac{2m+2}{2m+1}\cdot2P_m$\\
\midrule
$1$ & $2.6667$ & $3.5556$\\
$10$ & $3.0677$ & $3.2138$\\
$100$ & $3.1338$ & $3.1494$\\
$1000$ & $3.1408$ & $3.1424$\\
\bottomrule
\end{tabular}
\end{center}
$\pi=3.14159\ldots$は,どの$m$でも左右の値の間にある.有限の計算の結果が,どれほど確かな範囲を与えるかまで分かったのだ.ただし,区間の幅は$\frac{2P_m}{2m+1}$であり,\,$\frac{\pi}{2m+1}$より小さいものの,\,$m$に反比例する程度でしか縮まらない.\,$m=1000$でも,確定するのは$3.14$までである.
%
\subsection{二項係数}

偶数の公式の分数は,階乗で書き直せる.偶数の積からは$2$を1つずつくくり出せるので
\[
 (2m)!!=2^m m!
\]
である.また$(2m)!$は偶数の積と奇数の積に分かれるので
\[
 (2m-1)!!=\frac{(2m)!}{(2m)!!}=\frac{(2m)!}{2^m m!}
\]
となる.これらを式\eqref{式:ウォリス偶数}に代入すると
\begin{equation}\label{式:ウォリスと二項係数}
 I_{2m}=\frac{\pi}{2}\cdot\frac{(2m)!}{2^{2m}(m!)^2}
 =\frac{\pi}{2}\cdot\frac{1}{4^m}\binom{2m}{m}
\end{equation}
を得る.ここで$\binom{2m}{m}=\frac{(2m)!}{(m!)^2}$は二項係数であり,\,$2m$個から$m$個を選ぶ組合せの数${}_{2m}\mathrm{C}_m$に等しい.

二項係数$\binom{2m}{m}$は$m$とともに急速に大きくなる.では,\,$4^m$で割った値はどうなるだろうか.\,$\sqrt m$を掛けて調べてみよう.
\begin{center}
\begin{tabular}{cccccc}
\toprule
$m$ & $1$ & $10$ & $100$ & $1000$ & $10000$\\
\midrule
$\frac{\sqrt m}{4^m}\binom{2m}{m}$ & $0.5$ & $0.5572$ & $0.5635$ & $0.5641$ & $0.5642$\\
\bottomrule
\end{tabular}
\end{center}
値は$0.564$付近へ近づいているように見える.この値を,ウォリスの積分から求めよう.

偶数の公式と奇数の公式を掛けると,分子と分母の偶数がすべて約分されて
\[
 I_{2m}I_{2m+1}
 =\frac{1\cdot3\cdots(2m-1)}{3\cdot5\cdots(2m+1)}\cdot\frac{\pi}{2}
 =\frac{\pi}{2(2m+1)}
\]
となる.これと式\eqref{式:ウォリス隣接比}の$\frac{I_{2m+1}}{I_{2m}}\to1$を組み合わせれば
\begin{align*}
 mI_{2m}^2
 &=m\,I_{2m}I_{2m+1}\cdot\frac{I_{2m}}{I_{2m+1}}\\
 &=\frac{m\pi}{2(2m+1)}\cdot\frac{I_{2m}}{I_{2m+1}}
 \longrightarrow\frac{\pi}{4}
\end{align*}
となる.\,$I_{2m}>0$なので正の平方根を取ることができ,
\[
 \sqrt m\,I_{2m}\longrightarrow\frac{\sqrt\pi}{2}
\]
である.式\eqref{式:ウォリスと二項係数}より$\frac{\sqrt m}{4^m}\binom{2m}{m}=\frac2\pi\sqrt m\,I_{2m}$なので
\begin{equation}\label{式:中央二項係数の極限}
 \lim_{m\to\infty}\frac{\sqrt m}{4^m}\binom{2m}{m}
 =\frac1{\sqrt\pi}
\end{equation}
を得る.表の$0.564$付近の値は,\,$\frac1{\sqrt\pi}=0.56418\ldots$だったのだ.階乗を含む数列の大きさが,三角関数の積分を経由して分かったことになる.

\paragraphbf{硬貨を投げる確率}
硬貨を$2m$回投げて表がちょうど$m$回出る確率は,\,$\frac1{4^m}\binom{2m}{m}$である.式\eqref{式:中央二項係数の極限}によれば,この確率は$m$が大きいとき$\frac{1}{\sqrt{\pi m}}$に近い.例えば$200$回投げて表がちょうど$100$回出る確率は,\,$\frac1{\sqrt{100\pi}}=0.0564\ldots$に近く,実際の値は$0.0563\ldots$である.表と裏が半分ずつになることが最も起こりやすいとはいえ,ちょうど半分になる確率そのものは$0$へ近づいていく.
%
%
\section{ベータ積分}\label{節:ベータ積分}

\subsection{整数指数の場合}

0以上の整数$m,n$に対して
\[
 J_{m,n}=\int_0^1x^m(1-x)^n\dd x
\]
を考える.\,$m,n$が小さければ展開で求められる.しかし指数が大きくなると,展開した項数も多くなる.ここでは$x^m$と$(1-x)^n$をそのまま残し,部分積分で$(1-x)$の指数を1つずつ下げて,\,$n=0$の場合に帰着させることを目標とする.

$n\geqq1$として,\,$x^m$を積分し,\,$(1-x)^n$を微分すると
\begin{align*}
 J_{m,n}
 &=\left[\frac{x^{m+1}}{m+1}(1-x)^n\right]_0^1
   +\frac{n}{m+1}\int_0^1x^{m+1}(1-x)^{n-1}\dd x\\
 &=\frac{n}{m+1}J_{m+1,n-1}
\end{align*}
となる.\,$x=0$では$x^{m+1}$が,\,$x=1$では$n\geqq1$により$(1-x)^n$が$0$になるので,端点の項は消える.

添字が2つあるが,指数の和$m+n$は変わらず,\,$x$の指数が1つ増えて$(1-x)$の指数が1つ減っている.よって$n$が$0$になるまで繰り返せば
\begin{align*}
 J_{m,n}
 &=\frac{n}{m+1}J_{m+1,n-1}\\
 &=\frac{n}{m+1}\cdot\frac{n-1}{m+2}J_{m+2,n-2}\\
 &\qquad\vdots\\
 &=\frac{n}{m+1}\cdot\frac{n-1}{m+2}\cdots\frac{1}{m+n}J_{m+n,0}\\
 &=\frac{n!}{(m+1)(m+2)\cdots(m+n)}J_{m+n,0}
\end{align*}
となる.最後の積分は
\[
 J_{m+n,0}=\int_0^1x^{m+n}\dd x=\frac1{m+n+1}
\]
である.分母は$(m+1)(m+2)\cdots(m+n+1)$となり,分子と分母に$m!$を掛ければ分母は$(m+n+1)!$になるので
\begin{equation}\label{式:整数ベータ積分}
 \boxed{
 \int_0^1x^m(1-x)^n\dd x=\frac{m!\,n!}{(m+n+1)!}}
\end{equation}
を得る.\,$n=0$の場合も,右辺は$\frac{m!}{(m+1)!}=\frac1{m+1}$となり,同じ式で表される.

\begin{bookpoint}
\textbf{POINT}\par
$x^m$を積分し$(1-x)^n$を微分する部分積分で$J_{m,n}=\frac{n}{m+1}J_{m+1,n-1}$とし,\,$n=0$まで下げて$\frac{m!\,n!}{(m+n+1)!}$を得る.
\end{bookpoint}

\begin{example}{ベータ積分}
次の積分を求めよ.
\[
 \int_0^1x^3(1-x)^2\dd x.
\]
\end{example}

\paragraphbf{方針}
式\eqref{式:整数ベータ積分}で$m=3$,\,$n=2$とする.展開しても検算できる.

\paragraphbf{解答}
$\frac1{60}$

\paragraphbf{解法}
式\eqref{式:整数ベータ積分}で$m=3$,\,$n=2$とすると
\[
 \int_0^1x^3(1-x)^2\dd x
 =\frac{3!\,2!}{6!}
 =\frac{12}{720}
 =\frac1{60}
\]
となる.展開で検算すると
\begin{align*}
 \int_0^1(x^3-2x^4+x^5)\dd x
 &=\frac14-\frac25+\frac16\\
 &=\frac1{60}
\end{align*}
となり,一致する.公式があれば展開は要らないが,指数がどこに入るかを確かめるため,小さい場合を計算してみることにも意味がある.
%
\subsection{実数指数の場合}

指数を整数以外にも広げよう.\,$p,q>0$として
\[
 x^{p-1}(1-x)^{q-1}
\]
を考えたい.例えば$p=\frac12$なら$x^{p-1}=\frac1{\sqrt x}$となり,\,$x=0$で値が定まらない.閉区間で連続な関数の積分という,これまでの定義をそのまま使うことはできない.

そこで,問題のある端点をいったん避けて積分し,その端点へ近づけてみよう.\,$\frac1{\sqrt x}$と,比べるために$\frac1x$について,\,$\int_\varepsilon^1$を計算すると次のようになる.
\begin{center}
\begin{tabular}{ccccc}
\toprule
$\varepsilon$ & $0.1$ & $0.01$ & $0.0001$ & $0.000001$\\
\midrule
$\int_\varepsilon^1\frac1{\sqrt x}\dd x=2-2\sqrt\varepsilon$ & $1.3675$ & $1.8$ & $1.98$ & $1.998$\\[3pt]
$\int_\varepsilon^1\frac1x\dd x=-\log\varepsilon$ & $2.3026$ & $4.6052$ & $9.2103$ & $13.8155$\\
\bottomrule
\end{tabular}
\end{center}
どちらの関数も$x\to+0$で限りなく大きくなるが,\,$\frac1{\sqrt x}$の積分は$2$へ近づき,\,$\frac1x$の積分は限りなく大きくなる.関数の高さが限りなく大きくなることと,面積が限りなく大きくなることは同じではないのだ.

この考え方で積分を定める.\,$0<c<1$を1つ固定して
\begin{equation}\label{式:両端広義積分の定義}
 \int_0^1 f(x)\dd x
 =\lim_{\varepsilon\downarrow0}\int_{\varepsilon}^{c}f(x)\dd x
  +\lim_{\delta\downarrow0}\int_c^{1-\delta}f(x)\dd x
\end{equation}
と定める.ただし,右辺の\emph{2つの極限がそれぞれ有限の値を持つ場合}に限る.このような積分を\textbf{広義積分}といい,有限の値を持つことを収束するという.\,$c$を別の内点に変えても,その2点間の通常の積分が一方の極限から他方へ移るだけなので,和の値は変わらない.

両端へ同じ速さで近づける$\int_{\varepsilon}^{1-\varepsilon}$の極限だけでは,この定義を満たしたことにはならない.一方の端で正に発散し,他方の端で負に発散するものが打ち消し合うおそれがあるためである.この章のベータ積分は0以上の関数の積分だが,後にも使えるよう,両端を別々に調べる定義にしておく.

まず,\,$p>0$なら
\[
 \int_{\varepsilon}^{c}x^{p-1}\dd x
 =\frac{c^p-\varepsilon^p}{p}
 \longrightarrow\frac{c^p}{p}
\]
である.\,$p$が$1$より小さくても,\,$\varepsilon^p\to0$となるので収束する.表の$\frac1{\sqrt x}$は$p=\frac12$の場合である.

ベータ積分の左半分では$c=\frac12$とする.\,$x\in\left[0,\frac12\right]$において$(1-x)^{q-1}$は連続なので,ある正の定数$M$以下である.したがって
\[
 0\leqq x^{p-1}(1-x)^{q-1}\leqq Mx^{p-1}
\]
となる.下端$\varepsilon$を小さくすると積分値は増加するが,上の評価により$\frac{M}{p}\left(\frac12\right)^p$以下にとどまる.単調で上に有界な数列の極限と同じ考え方で,この極限は有限の値として存在する.右半分も$u=1-x$と置けば,\,$q>0$によって同様に収束する.したがって$p,q>0$なら,両端の広義積分がともに定まる.

\begin{definition}{ベータ関数}
$p,q>0$に対して
\begin{equation}\label{式:ベータ関数の定義}
 B(p,q)=\int_0^1x^{p-1}(1-x)^{q-1}\dd x
\end{equation}
と定める.端点で関数が定まらない場合には,式\eqref{式:両端広義積分の定義}による広義積分とする.
\end{definition}

整数の場合との対応は,指数が$1$だけずれている.式\eqref{式:整数ベータ積分}は
\[
 B(m+1,n+1)=\frac{m!\,n!}{(m+n+1)!}
\]
と書ける.\,$B(m,n)$の被積分関数の指数は$m-1,n-1$であり,\,$m,n$そのものではない.記号の引数と被積分関数の指数を取り違えないようにする.
%
\subsection{対称性と漸化式}

$u=1-x$と置くと2つの指数が入れ替わるため
\begin{equation}\label{式:ベータ対称性}
 B(p,q)=B(q,p)
\end{equation}
である.広義積分の場合も,端点を避けた区間で置換し,各端点の極限を取れば同じ結論になる.

次に,整数の場合と同じく,指数を1つずらした積分どうしの関係を作ろう.被積分関数に$x+(1-x)=1$を掛ければ
\begin{equation}\label{式:ベータ和の関係}
 B(p+1,q)+B(p,q+1)=B(p,q)
\end{equation}
となる.もう1つの関係を得るために$x^p(1-x)^q$を微分する.
\[
 \frac{\dd}{\dd x}\{x^p(1-x)^q\}
 =p x^{p-1}(1-x)^q-qx^p(1-x)^{q-1}
\]
これを$\varepsilon$から$1-\delta$まで積分し,両端をそれぞれ極限へ近づける.\,$p,q>0$なので$x^p(1-x)^q$は両端で$0$へ近づき,右辺の2つの積分はすでに収束すると分かっている.よって
\[
 0=pB(p,q+1)-qB(p+1,q)
\]
を得る.これは整数の場合の部分積分を,端点を避けて行ったものである.式\eqref{式:ベータ和の関係}に$B(p,q+1)=\frac qpB(p+1,q)$を代入すると
\[
 \frac{p+q}{p}B(p+1,q)=B(p,q)
\]
となるので
\begin{equation}\label{式:ベータ漸化式}
 B(p+1,q)=\frac{p}{p+q}B(p,q),\qquad
 B(p,q+1)=\frac{q}{p+q}B(p,q)
\end{equation}
を得る.この漸化式は整数以外にも成り立つ.

整数のときには,\,$B(1,1)=\int_0^1 1\dd x=1$から出発して階乗の公式を再び導ける.右の式を$q=1,2,\ldots,n$について順に使うと
\[
 B(1,n+1)=\frac{n}{n+1}\cdot\frac{n-1}{n}\cdots\frac12B(1,1)=\frac{1}{n+1}
\]
となる.次に左の式を$p=1,2,\ldots,m$について順に使うと
\begin{align*}
 B(m+1,n+1)
 &=\frac{m}{m+n+1}\cdot\frac{m-1}{m+n}\cdots\frac{1}{n+2}B(1,n+1)\\
 &=\frac{m!\,n!}{(m+n+1)!}
\end{align*}
となり,式\eqref{式:整数ベータ積分}と一致する.どちらも比が添字の式で与えられる階比型であり,比を順に掛けて一般項を得ている.
%
\subsection{三角関数による表示}

$x^{p-1}(1-x)^{q-1}$の$x$と$1-x$は,和が$1$である.和が$1$になる2つの量といえば$\sin^2t$と$\cos^2t$であろう.そこで$x=\sin^2t$と置くと,\,$1-x=\cos^2t$であり,\,$0<t<\frac{\pi}{2}$で
\[
 \dd x=2\sin t\cos t\dd t
\]
である.この範囲では$\sin t,\cos t$が正なので,実数の指数もそのまま扱える.式\eqref{式:ベータ関数の定義}は
\begin{equation}\label{式:ベータ三角表示}
 B(p,q)=2\int_0^{\frac{\pi}{2}}
 (\sin t)^{2p-1}(\cos t)^{2q-1}\dd t
\end{equation}
となる.端点で関数が定まらない場合は,端点を避けて置換した後で極限を取る.左辺の両端の積分が収束しているので,対応する右辺の極限も有限である.

$p=q=\frac12$なら両方の指数が$0$になり,
\[
 B\left(\frac12,\frac12\right)
 =2\int_0^{\frac{\pi}{2}}1\dd t=\pi
\]
である.つまり
\[
 \int_0^1\frac1{\sqrt{x(1-x)}}\dd x=\pi
\]
となる.両端で高さが限りなく大きくなるにもかかわらず,積分は有限になる具体例である.

一方,\,$p=\frac{n+1}{2}$,\,$q=\frac12$とすれば,\,$\cos$の指数が$0$になって
\[
 B\left(\frac{n+1}{2},\frac12\right)=2\int_0^{\frac\pi2}\sin^nt\dd t=2I_n
\]
となる.\,$p$を$1$増やすことは,\,$\sin$の指数$n$を$2$増やすことに対応する.実際,式\eqref{式:ベータ漸化式}で$p=\frac{n+1}2$,\,$q=\frac12$とすると,\,$\frac{p}{p+q}=\frac{n+1}{n+2}$なので
\[
 B\left(\frac{n+3}{2},\frac12\right)
 =\frac{n+1}{n+2}
  B\left(\frac{n+1}{2},\frac12\right)
\]
となる.両辺を$2$で割れば$I_{n+2}=\frac{n+1}{n+2}I_n$であり,ウォリスの漸化式\eqref{式:ウォリス漸化式}そのものである.整数指数の多項式の積分とウォリスの積分は,置換$x=\sin^2t$によって,同じベータ積分の異なる場合として結び付いたのだ.
%
%
\section*{章末問題}
\begin{enumerate}[label=\textbf{\thechapter.\arabic*},leftmargin=*]
 \item\label{問:区分求積分子が添字}
 次の極限を求めよ.
 \[
  \lim_{n\to\infty}\sum_{k=1}^{n}\frac{k}{n^2+k^2}.
 \]
 \item\label{問:区分求積根号}
 次の極限を求めよ.
 \[
  \lim_{n\to\infty}\cfrac1{n^{\cfrac32}}
      \sum_{k=1}^{n}\sqrt{n+k}.
 \]
 \item\label{問:逆平方和の残り}
 自然数$N>n\geqq1$に対して
 \[
  \frac1{n+1}-\frac1{N+1}
  \leqq\sum_{k=n+1}^{N}\frac1{k^2}
  \leqq\frac1n-\frac1N
 \]
 を示せ.これを使って$\sum_{k=1}^{\infty}\frac1{k^2}$が収束し,その第$n$項より後の和$R_n$が
 $\frac1{n+1}\leqq R_n\leqq\frac1n$を満たすことを示せ.
 \item\label{問:積分制約と二乗}
 $f$は$[0,1]$で連続で,\,$\int_0^1xf(x)\dd x=1$を満たすとする.
 $\int_0^1f(x)^2\dd x$の最小値を求め,それを与える関数を求めよ.
 \item\label{問:対数の積分漸化式}
 非負整数$n$に対して$K_n=\int_1^{\e}(\log x)^n\dd x$とおく.
 $K_0$と,\,$n\geqq1$における漸化式を求め,\,$K_3$を計算せよ.
 \item\label{問:ウォリス積の誤差}
 $P_m=\prod_{k=1}^{m}\frac{4k^2}{4k^2-1}$とする.
 \[
 0\leqq\frac{\pi}{2}-P_m\leqq\frac{\pi}{4m+4}
 \]
 を示せ.また$\lim_{m\to\infty}\sqrt m\int_0^{\frac{\pi}{2}}\cos^{2m}x\dd x$を求めよ.
 \item\label{問:整数ベータ二項係数}
 非負整数$m,n$に対して
 \[
 \int_0^1\binom{m+n}{m}x^m(1-x)^n\dd x=\frac1{m+n+1}
 \]
 を示せ.次に$\int_0^1x^2(1-x)^4\dd x$を求めよ.
 \item\label{問:半整数ベータ}
 広義積分$\int_0^1\sqrt{\frac{x}{1-x}}\dd x$の収束を両端で確かめ,ベータ関数の漸化式を使って値を求めよ.さらに$x=\sin^2t$による置換でも検算せよ.
\end{enumerate}

\smallskip\noindent 解答は\pageref{解答:第10章}ページ.

%
